% ============================================================================
%  B — file phần: Đọc sâu
%  Ngày tạo: 2026-09-06 | Sửa gần nhất: 2026-09-06 | Ôn gần nhất: chưa
% ============================================================================
\documentclass[../english.tex]{subfiles}
\begin{document}
\part{ĐỌC SÂU}
\begin{tomtat}
Phần B là phần trung tâm với mục tiêu của quyển sách: lấy đúng và đủ ý từ một văn bản khó. Sáu chương chia theo \emph{việc bạn đang làm khi đọc}: chọn tầng đọc phù hợp với mục đích (5); nhận ra bộ khung lập luận của tác giả (6); xử lý một bài báo khoa học theo ba lượt (7); đọc ra những gì tác giả ngụ ý mà không nói thẳng (8); theo kịp một bài giảng hoặc khoá học (9); và xây vốn từ chuyên ngành từ chính tài liệu bạn đọc (10). Nếu chỉ nhớ một điều: người đọc giỏi không đọc nhanh hơn --- họ \emph{quyết định trước} lần đọc này nhằm mục đích gì, và mỗi mục đích có một cách đọc khác nhau.
\end{tomtat}

Có một điều mà người đọc tài liệu chuyên ngành hiếm khi được dạy: phần lớn thời gian đọc bị lãng phí không phải vì đọc chậm mà vì \emph{đọc sai tầng}. Người ta đọc kỹ từng câu một bài báo mà lẽ ra chỉ cần lướt để biết nó có liên quan không; hoặc ngược lại, lướt qua một bài quan trọng rồi tưởng đã hiểu và trích dẫn sai ý tác giả. Chương 5 xử lý đúng vấn đề đó.

Ba chương tiếp theo đi vào ba loại thông tin mà một văn bản học thuật mang, và chúng nằm ở ba chỗ khác nhau. Nội dung \emph{mệnh đề} --- tác giả khẳng định gì --- nằm rải trong các câu. Bộ khung \emph{lập luận} --- các khẳng định nối với nhau ra sao, cái nào là bằng chứng cho cái nào --- nằm ở các từ nối và ở cấu trúc đoạn, và đó là chương 6. Còn \emph{thái độ và mức cam kết} --- tác giả tin điều này tới đâu, họ ngầm phê bình ai --- nằm ở những từ nhỏ mà người học hay lướt qua, và đó là chương 8, chương quan trọng nhất phần này cho mục tiêu ``hiểu đúng ý người nói''.

Chương 7 và 9 áp ba tầng đó vào hai loại tài liệu cụ thể mà bạn dùng hằng ngày: bài báo và bài giảng. Chương 10 khép lại bằng việc biến những gì bạn đọc thành vốn từ ngành của riêng bạn --- vì như \ptr{A/02} đã chỉ ra, đó là con đường ngắn nhất tới ngưỡng đọc trôi chảy.
\subfile{00-map}
\subfile{05-ba-tang-doc/ba-tang-doc}
\subfile{06-cau-truc-lap-luan/cau-truc-lap-luan}
\subfile{07-doc-bai-bao-khoa-hoc/doc-bai-bao-khoa-hoc}
\subfile{08-doc-giua-cac-dong/doc-giua-cac-dong}
\subfile{09-nghe-giang/nghe-giang}
\subfile{10-tu-vung-chuyen-nganh/tu-vung-chuyen-nganh}
\ifSubfilesClassLoaded{\printbibliography[title={Nguồn}]}{}
\end{document}
